\section{Common-divisor strata and intrinsic families}
\label{sec:common-divisor-strata}

The coefficient condition in Theorem~\ref{thm:recognition-dichotomy-v148}
is not imposed in this section.  We instead retain the common divisor
of a space of first relations.  This geometric condition gives a
larger smooth parameter space, includes relations without either
factor symmetry, and has a scheme-theoretic inverse.  The distinction
between a stratum and the condition on its geometric points is
essential when the parameter base is nonreduced.

Write \(S_j(E)=\Sym^j E\), where \(E\) is a complex vector space.
For \(r\ge2\), denote by
\[
 X_{h,r}(E)\subset\Gr(r,S_h(E))
\]
the locus of subspaces with greatest common divisor equal to \(1\),
up to a scalar.  Empty Grassmannians are omitted throughout.  When
\(h=0\) this notation can also be used for \(r=1\).

\begin{lemma}[Universal common divisor]\label{lem:universal-gcd}
Let \(D=g+h\).  The spaces of \(r\) degree-\(D\) forms whose exact
greatest common divisor has degree \(g\) form a locally closed
reduced subscheme \(Z_g\) of \(\Gr(r,S_D(E))\).  Multiplication
induces an isomorphism
\begin{equation}\label{eq:universal-gcd}
 \Pj(S_g(E))\times X_{h,r}(E)\ \xrightarrow{\ \sim\ }\ Z_g,
                 \qquad ([f],J)\longmapsto fJ.
\end{equation}
In particular, the common-divisor line and the residual subbundle of
a family with classifying map in \(Z_g\) are recovered uniquely and
commute with arbitrary complex base change.
\end{lemma}
\begin{proof}
For each \(a\), multiplication from
\(\Pj(S_a(E))\times\Gr(r,S_{D-a}(E))\) has closed image: its domain
is projective.  The union for \(a>g\) is closed.  Remove this union
from the image for \(a=g\), and give the resulting locally closed
set its reduced structure.  Its points are exactly those in the
statement.  The inverse image of this set under multiplication is
\(\Pj(S_g(E))\times X_{h,r}(E)\).  The same closed-image argument
with \(D=h\) shows that \(X_{h,r}(E)\) is open.

The restricted multiplication map is proper, by base change.  It
is injective on geometric points over every algebraically closed
extension of \(\C\), since a greatest common divisor is unique up
to scalar in a polynomial ring.  Its differential is injective as
well.  Indeed, a tangent vector in its kernel satisfies
\[
                   \dot f\, j+f\dot j\in fJ\quad(j\in J).
\]
Consequently \(f\mid\dot f\,j\) for every \(j\in J\).  For each
irreducible factor of \(f\), some element of \(J\) is not divisible
by that factor.  Comparison of its valuation gives
\(f\mid\dot f\).  Equality of degrees makes \(\dot f\) a scalar
multiple of \(f\); then \(\dot J=0\) in
\(\Hom(J,S_h(E)/J)\).

In characteristic zero the geometric-point injectivity and this
differential calculation imply that the map is universally
injective and unramified.  A proper, universally injective,
unramified morphism is a closed immersion
\cite[Lemma~41.7.2]{StacksClosedImmersion}.
Here its image is all of the reduced target.  The defining ideal
of a closed immersion with this support is zero in a reduced
scheme.  This proves the isomorphism, including its assertion for
nonreduced test schemes.  Equivalently, one can check properness
on the base change of the original projective multiplication map;
that base change has the same underlying open preimage.  Its source
is reduced because it is this open subscheme of the original
smooth product.
\end{proof}

\begin{remark}\label{rem:gcd-scheme-condition}
A map from a nonreduced scheme is required to factor through the
locally closed \emph{scheme} \(Z_g\), not merely to have geometric
points in its underlying set.  These are different requirements.
For example, over \(\C[\epsilon]/(\epsilon^2)\) the span of
\(x^2+\epsilon y^2,xy\) specializes to a pair with common divisor
\(x\), but has no factor \(x+\epsilon(ax+by)\).  Comparison of the
\(y^2\) coefficient in the first form proves this.  The stratum in
Lemma~\ref{lem:universal-gcd} excludes this transverse infinitesimal
family.  No use of a fibrewise radical can replace its scheme structure.
\end{remark}

Now let \(\dim V=\dim U=n\ge3\),
\(E=V^*\otimes U\), and \(\delta=\det T\).  Denote by
\(G\subset\operatorname{GL}(E)\) the stabilizer of the line
\(\C\delta\).  By Lemma~\ref{lem:linear-preserver-v147},
\begin{equation}\label{eq:determinant-group}
 G^\circ=(\operatorname{GL}(V)\times\operatorname{GL}(U))/\mathbb G_m,
                 \qquad G/G^\circ=\mathbf Z/2\mathbf Z.
\end{equation}
The second component exchanges the two tensor factors; the diagonal
scalars in the displayed product act trivially on the matrix space.
Fix \(s\ge1\) and \(D=sn+h\).  Inside \(Z_{sn}\), impose the condition
that the common-divisor line is the \(s\)-th power of a determinant
form after a linear change of the \(n^2\) variables.  Call the resulting
locally closed scheme \(Z^{\det}_{s,h,r}\).

\begin{theorem}[Intrinsic determinant-divisor stratum]
\label{thm:determinant-stratum}
There are canonical isomorphisms of schemes and algebraic stacks
\begin{equation}\label{eq:determinant-stratum}
 Z^{\det}_{s,h,r}\simeq
 \operatorname{GL}(E)\times^G X_{h,r}(E),\qquad
 [Z^{\det}_{s,h,r}/\operatorname{GL}(E)]\simeq[X_{h,r}(E)/G].
\end{equation}
They hold over arbitrary complex parameter schemes.  The stack is
smooth.  At a residual relation \(J\), its tangent complex, in degrees
\(-1,0\), is
\begin{equation}\label{eq:ambient-tangent-complex}
 [\,\operatorname{Lie}(G)\longrightarrow
                    \Hom(J,S_h(E)/J)\,],
             \qquad X\longmapsto(j\mapsto Xj\bmod J).
\end{equation}
No invariance of \(J\) under either full factor group is assumed.
The equivalence identifies stabilizers, infinitesimal families, and
specializations whose classifying maps remain in this stratum.
\end{theorem}
\begin{proof}
The power map \(\Pj(S_n(E))\to\Pj(S_{sn}(E))\),
\([f]\mapsto[f^s]\), is a closed immersion in characteristic zero.
It is proper and geometrically injective, and its differential has
kernel given by \(sf^{s-1}\dot f\in\C f^s\), hence by
\(\dot f\in\C f\).  The same closed-immersion criterion as above
applies.  The orbit of \([\delta]\) is the smooth locally closed
homogeneous space \(\operatorname{GL}(E)/G\).  The preimage of its
power orbit under the first projection in
\eqref{eq:universal-gcd} is therefore a locally closed subscheme.
It is exactly \(Z^{\det}_{s,h,r}\).

Pulling back along the \(G\)-torsor
\(\operatorname{GL}(E)\to\operatorname{GL}(E)/G\) trivializes the
common-divisor form, leaving the unrestricted residual relation
\(J\in X_{h,r}(E)\).  Descent gives the first isomorphism in
\eqref{eq:determinant-stratum}; descent of equivariant maps and
morphisms gives the second.  This proof uses morphisms of schemes,
not a comparison only on complex points.  It also supplies local
lifts of all isomorphisms over an fppf cover.

The space \(X_{h,r}(E)\) is open in a Grassmannian and \(G\) is
smooth.  The quotient is algebraic and smooth
\cite[Theorem~97.17.2]{StacksQuotient}; see also
\cite[Lemma~101.33.7]{StacksSmoothQuotient}.
For completeness, write a first-order residual subspace as the
graph of \(\epsilon\phi:J\to S_h(E)/J\).  An infinitesimal change
\(1+\epsilon X\) replaces \(\phi\) by \(\phi+X|_J\bmod J\).
This proves the tangent complex.  Its kernel records infinitesimal
automorphisms and its cokernel records first-order isomorphism
classes.  The same torsor description proves compatibility with
specialization.  No assertion is made about limits outside the
specified common-divisor stratum.
\end{proof}

The next step connects these parameter spaces with ungraded finite
algebras.  It is useful to state the family condition intrinsically,
since an absolute nilradical is unsuitable over a nonreduced base.
Put \(e=\dim E\) and
\(\ell=\binom{e+D}{D}-r\).  Consider finite locally free commutative
unital algebras \(\mathcal A\) of rank \(\ell\) on a complex scheme
\(T_0\), subject to the following conditions.  The map
\(\epsilon(a)=\ell^{-1}\operatorname{tr}(M_a)\) is multiplicative.
For \(\mathcal N=\ker\epsilon\), one has
\(\mathcal N^{D+1}=0\); all quotients of its power filtration are
locally free; \(\mathcal E=\mathcal N/\mathcal N^2\) has rank \(e\);
and multiplication gives isomorphisms
\(\Sym^j\mathcal E\simeq\mathcal N^j/\mathcal N^{j+1}\) for
\(j<D\).  In degree \(D\) its kernel has rank \(r\).  These
conditions specify the maximal-lower-Hilbert-function stratum; they
do not assert that arbitrary smoothings of a local algebra remain
in it.

\begin{proposition}[Relative first-relation algebra]
\label{prop:relative-first-relation}
Every such algebra is locally isomorphic, over its base, to
\[
 \mathcal A(\mathcal K)=
 \Sym\mathcal E/((\mathcal K)+\mathfrak m^{D+1}),\qquad
 \mathcal K=\ker(\Sym^D\mathcal E\longrightarrow\mathcal N^D).
\]
Its trace augmentation, \(\mathcal E\), and \(\mathcal K\) are
intrinsic and compatible with arbitrary base change in this
stratum.  For two such algebras, taking the linear part gives a
locally surjective map of isomorphism sheaves to the sheaf of
isomorphisms of \((\mathcal E,\mathcal K)\).  Its kernel is the
smooth substitution group whose underlying scheme in a local
presentation is
\(\Hom(\mathcal E,\mathcal N^2)\).

After quotienting morphisms by this tangent-identity kernel and
fppf stackifying, the algebras with first-relation map in
\(Z^{\det}_{s,h,r}\) have stack \([X_{h,r}(E)/G]\).
\end{proposition}
\begin{proof}
Locally lift a basis of \(\mathcal E\) to \(\mathcal N\).
Nilpotence and multiplication on the successive quotients make the
induced map from the truncated symmetric algebra surjective.
A relation with a nonzero term of least degree \(j<D\) would give
a nonzero kernel in degree \(j\) of the associated graded algebra,
which is impossible.  Its entire kernel therefore lies in degree
\(D\) and is precisely \(\mathcal K\).  Changing a lift by an
element of \(\mathcal N^2\) does not alter a product of \(D\)
generators, since the difference has degree at least \(D+1\).
This also proves the asserted uniqueness of the first relation.

In this presentation multiplication by a positive-degree element
is strictly filtration-increasing.  Its trace is zero over any
base ring.  Multiplication by a scalar has trace \(\ell\) times
that scalar, proving that the presentation's augmentation is the
specified intrinsic trace augmentation.  All filtration sequences
are locally split sequences of vector bundles, so their kernels
and quotients commute with base change.

An isomorphism must preserve trace and induces a linear map carrying
\(\mathcal K\) to \(\mathcal K'\).  Conversely such a map lifts
locally to an isomorphism by the homogeneous presentations.  All
other lifts are obtained by substituting generators plus arbitrary
elements of \(\mathcal N'^2\); the inverse is constructed degree by
degree in a finite filtration.  This gives the stated kernel and
its smoothness, without replacing its substitution law by addition.
After identifying morphisms with the same linear part, frames of
\(\mathcal E\) give the quotient
\([Z^{\det}_{s,h,r}/\operatorname{GL}(E)]\).
Theorem~\ref{thm:determinant-stratum} completes the proof.
\end{proof}

\begin{remark}\label{rem:algebra-inertia}
The original algebra stack is not this quotient stack: it has the
additional nonlinear inertia just computed.  The phrase
\emph{effective first-relation stack} will refer to the specified
quotient on morphisms followed by stackification.  The construction
retains the full automorphism group before taking that quotient.
It does not identify algebras with the same geometric fibres, nor
does it discard nilpotent parameters.  In particular, its relation
to a stratum of first relations is stronger than a bijection of
isomorphism classes over \(\C\).
\end{remark}
